LLM evaluation has progressed from broad tests of knowledge and reasoning
\citep{hendrycks2021,srivastava2023,liang2023} to benchmarks built around
professional and economically valuable work
\citep{gdpval2025,profbench2025,occubench2026,jobbench2026}. These benchmarks
measure an important capability: whether a model can produce a strong
deliverable on its own. In deployed workflows, however, a model may instead
assist a human or another model by planning, decomposing, checking, or revising
while the downstream worker produces the final output. Standalone performance
need not predict performance in such interactive settings
\citep{collins2024,chang2025}. The model that is the best direct solver may not
be the best assistant.

Existing work does not provide a general cross-model comparison of these two
roles. Human--AI field studies establish that assistance can change
productivity, but typically evaluate one tool in one domain
\citep{noy2023,peng2023,dellacqua2023,brynjolfsson2025}. Multi-agent and
tutoring studies examine teaching, planning, or complementarity in narrower
settings \citep{saha2023,bansal2021,macina2025,gandhi2025}. These studies do
not rank the same models as both direct solvers and assistants on the same
tasks. This leaves unanswered a basic evaluation question: does a model's
automation rank predict the value of its guidance to a fixed downstream
worker?

We introduce \textsc{CentaurBench}, a controlled benchmark for measuring that
relationship. We evaluate nine focal LLMs on seven professional tasks under
two modes. In \emph{automation}, the focal model produces the deliverable
directly. In \emph{augmentation}, it produces process guidance for a fixed
GPT-3.5-Turbo worker, which then produces the deliverable. Holding the worker,
task, and worker prompt fixed isolates variation due to the assistant's
guidance rather than the executor. Outputs are scored through blinded,
rubric-guided pairwise comparisons by four LLM judges, with leave-family-out
masking to prevent a judge from evaluating outputs from its own provider
family. We repeat the complete generation and evaluation pipeline ten times
and report task-level and aggregate ranks.

The two modes yield materially different model profiles. Across the nine
models evaluated in both roles, average automation and augmentation ranks have
only moderate correlation ($\rho=0.48$), with task-level correlations ranging
from $-0.04$ to $0.85$. The winning model differs across modes on five of seven
tasks. Assistance is also not uniformly beneficial: the unaided worker ranks
first on three tasks, and only one assisted condition beats it on average.
Large task-level reversals persist across repeated runs; for example,
Claude-Opus-4.8 leads automation on market-trends analysis but is among the
weakest assistants on that task.

Our contribution is threefold. First, we operationalize role-aware model
evaluation as a cross-model, cross-task benchmark. Second, the fixed-worker
design provides a controlled measure of marginal guidance value. Third, the
resulting rank matrix shows that assistant quality is task-dependent and not
recoverable from an automation leaderboard alone. \textsc{CentaurBench}
therefore complements autonomous benchmarks with an evaluation axis directly
relevant to human-in-the-loop and multi-agent systems.
